% Full finite-framework proofs for P18-A1.
\section{Finite abstract argumentation: proofs}
\label{sec:annex-p18-finite-argumentation}

Use the definitions of Chapter~\ref{sec:p18-finite-argumentation}. These proofs
supply the stated finite instances of Dung's framework results and the
additional finite examples \cite[pp.~326--330]{Dung1995}.

\begin{annexitem}{annex:p18.defense-properties}{thm:p18-defense-properties}{Defense and admissible enlargement}
\begin{proof}
If \(S\subseteq T\), then \(S^+\subseteq T^+\): any attack with
source in \(S\) also has source in \(T\). Hence
\(\operatorname{Att}(a)\subseteq S^+\) implies
\(\operatorname{Att}(a)\subseteq T^+\), proving monotonicity.

Suppose \(S\) is conflict-free. If \(a,b\in\mathsf F(S)\) and
\((a,b)\in R\), defense of \(b\) supplies \(s\in S\) with
\((s,a)\in R\). Defense of \(a\) then supplies \(t\in S\) with
\((t,s)\in R\). This contradicts conflict-freeness of \(S\).
The argument includes \(a=b\), so it also excludes a self-attack in
\(\mathsf F(S)\). Thus \(\mathsf F(S)\) is conflict-free.

If \(S\) is admissible, then \(S\subseteq\mathsf F(S)\).
Monotonicity gives \(\mathsf F(S)\subseteq\mathsf F(\mathsf F(S))\),
and the preceding paragraph gives conflict-freeness of \(\mathsf F(S)\).
These are precisely its admissibility conditions.

Finally let \(a\in\mathsf F(S)\), with \(S\) admissible, and put
\(T=S\cup\{a\}\). Both \(S\) and \(a\) lie in the conflict-free
set \(\mathsf F(S)\), so \(T\) is conflict-free. Since
\(S\subseteq T\), monotonicity gives
\[
 T\subseteq\mathsf F(S)\subseteq\mathsf F(T).
\]
Thus \(T\) is admissible. The same inclusion
\(\mathsf F(S)\subseteq\mathsf F(T)\) says that every argument
previously defended remains defended. This proves both parts of the
admissible-enlargement assertion.
\end{proof}
\end{annexitem}

\begin{annexitem}{annex:p18.grounded-iteration}{thm:p18-grounded-iteration}{Grounded iteration and leastness}
\begin{proof}
The empty set is conflict-free and is contained in \(\mathsf F(\varnothing)\),
so it is admissible. The preceding theorem now proves by induction that
every \(G_i\) is admissible. Its admissibility gives
\(G_i\subseteq\mathsf F(G_i)=G_{i+1}\).

Every strict inclusion in this increasing sequence adds an argument of the
finite set \(A\). There can be at most \(n=|A|\) strict inclusions.
If there were no equality for \(0\leq i\leq n\), the first \(n+1\)
inclusions would all be strict, requiring more than \(n\) arguments.
Therefore \(G_k=G_{k+1}\) for some \(k\leq n\). It is a fixed
point, and subsequent stages equal it.

Let \(E\subseteq A\) be any fixed point of \(\mathsf F\), with no
conflict-freeness assumption. The base inclusion \(G_0\subseteq E\)
holds because \(G_0\) is empty. If \(G_i\subseteq E\), monotonicity
gives \(G_{i+1}\subseteq\mathsf F(E)=E\). Thus \(G_k\subseteq E\)
for every fixed point \(E\). This proves leastness. Two least fixed
points would each contain the other, so this one is unique.

By definition, a complete set is conflict-free, satisfies
\(S\subseteq\mathsf F(S)\) by admissibility, and satisfies the reverse
inclusion by completeness. Hence it is a conflict-free fixed point.
Conversely a conflict-free fixed point satisfies both inclusions and is
complete. The set \(G_k\) is admissible and fixed, so it is complete.
Leastness places it in every complete extension. Since it is itself one of
those extensions, their intersection is exactly \(G_k\).

For the stated procedure, each pass calculates \(S^+\) from the given
finite attack relation, then calculates exactly \(\mathsf F(S)\) by the
defense definition. The successive values of \(S\) are therefore
\(G_0,G_1,\ldots\). The equality is detected on evaluating
\(\mathsf F(G_k)\), after \(k+1\leq n+1\) evaluations. All scans
within an evaluation are finite. If \(A=\varnothing\), the first
evaluation returns the empty set and the equality test succeeds immediately.
\end{proof}
\end{annexitem}

\begin{annexitem}{annex:p18.extension-relations}{thm:p18-extension-relations}{Preferred existence and stable inclusion}
\begin{proof}
Fix an admissible set \(S\). The family of its admissible supersets is
nonempty, since it contains \(S\), and finite, since it is a subfamily
of \(2^A\). Choose a member \(E\) of largest cardinality in this
family. A strictly larger admissible superset of \(E\) would still
contain \(S\) and would have larger cardinality, a contradiction.
Thus \(E\) is a preferred extension containing \(S\). Cardinality is
used here only to select a maximal element of a finite family; it does not
change the inclusion-based definition of preferred extension. Taking
\(S=\varnothing\) proves existence.

Let \(E\) be preferred. If some \(a\in\mathsf F(E)\) were absent
from \(E\), admissible enlargement would give the strictly larger
admissible set \(E\cup\{a\}\). Maximality excludes this. Consequently
\(\mathsf F(E)\subseteq E\), and admissibility supplies the reverse
inclusion and conflict-freeness. Therefore \(E\) is complete.

Let \(E\) be stable, and take \(a\in E\) and an attacker \(b\)
of \(a\). Conflict-freeness forces \(b\notin E\). Stability therefore
supplies a member of \(E\) attacking \(b\), so \(E\) defends every
member and is admissible. If a strictly larger set \(T\) contained
\(E\), an argument \(b\in T\setminus E\) would be attacked by a
member of \(E\subseteq T\). Thus \(T\) would not be conflict-free,
and in particular would not be admissible. This proves that \(E\) is
preferred. The family inclusions follow, and the grounded containment is
the leastness already proved.

If \(S\) is stable, conflict-freeness gives
\(S\subseteq A\setminus S^+\), and attack on every outsider gives
\(A\setminus S^+\subseteq S\). Hence equality holds.
Conversely, if \(S=A\setminus S^+\), then \(S\cap S^+=\varnothing\),
so \(S\) is conflict-free, and each argument outside \(S\) lies in
\(S^+\). Thus \(S\) is stable. The complete finite calculations in
Annex~\ref{annex:p18.example-extensions} provide the asserted failures of
existence and of converse inclusion.
\end{proof}
\end{annexitem}

\begin{annexitem}{annex:p18.example-extensions}{thm:p18-example-extensions}{Complete calculations for the finite examples}
\begin{proof}
\paragraph{Chain.}
The conflict-free sets are exactly
\(\varnothing,\{a\},\{b\},\{c\},\{a,c\}\): any other subset
contains one of the two attacks. The argument \(b\) cannot be defended
by any set, because its attacker \(a\) has no attacker. The argument
\(a\) is always defended, while \(c\) is defended exactly when
the defending set contains \(a\), the sole attacker of \(b\).
Thus the admissible sets are exactly
\(\varnothing,\{a\},\{a,c\}\). On these sets the characteristic
operator has values \(\{a\},\{a,c\},\{a,c\}\), respectively.
This proves the listed grounded sequence and the sole complete/preferred
extension. The set \(\{a,c\}\) attacks the only outsider \(b\), so
it is stable; no other stable set is possible because stable sets are preferred.

\paragraph{Mutual pair.}
The conflict-free sets are precisely \(\varnothing,\{a\},\{b\}\).
Each singleton attacks the sole attacker of its member, so each is admissible.
No argument is unattacked, giving \(\mathsf F(\varnothing)=\varnothing\).
Also \(\mathsf F(\{a\})=\{a\}\) and
\(\mathsf F(\{b\})=\{b\}\): each singleton defends its member but
does not attack that member, so it cannot defend the opposite argument.
All three sets are complete, the singletons are the maximal admissible sets,
and each singleton attacks its one outsider. The empty set is not stable.

\paragraph{Self-attack.}
Only the empty set is conflict-free, so it is the only possible admissible,
complete or preferred set. It defends no argument, since \(a\) attacks
itself and the empty set attacks nothing. Hence
\(\mathsf F(\varnothing)=\varnothing\), proving the grounded and
complete claims. The empty set fails to attack its outsider \(a\), so
there is no stable extension. On the other hand, \(\{a\}\) attacks
the sole attacker of \(a\), namely \(a\), so
\(\mathsf F(\{a\})=\{a\}\). This fixed point is not conflict-free.

\paragraph{Empty framework.}
Its only subset is \(\varnothing\). It is conflict-free and defends all
its members vacuously. Its characteristic value is empty because there is
no argument to include. Hence it is grounded, complete and preferred. There
is no outsider to attack, so the same set is stable. A family containing this
one extension is \(\{\varnothing\}\), not the empty family.

\paragraph{Unequal sizes.}
The conflict-free sets are exactly
\(\varnothing,\{a\},\{b\},\{c\},\{b,c\}\). Every nonempty set
on this list defends its members: \(a\) attacks both its attackers, and
each of \(b,c\) attacks its sole attacker \(a\). Thus all listed
sets are admissible. There are no unattacked arguments, so
\(\mathsf F(\varnothing)=\varnothing\). Directly from the attacker
sets,
\[
 \mathsf F(\{a\})=\{a\},\qquad
 \mathsf F(\{b\})=\mathsf F(\{c\})
   =\mathsf F(\{b,c\})=\{b,c\}.
\]
This gives exactly the three complete sets in the statement and grounded
set \(\varnothing\). The maximal admissible sets are \(\{a\}\) and
\(\{b,c\}\); each attacks every outsider. Their sizes are one and two.

\paragraph{Shared argument.}
Any conflict-free set containing \(c\) must exclude \(a,b,d\), because
\(a,b\) attack \(c\) and \(c\) attacks \(d\). But \(\{c\}\)
does not attack either of its attackers, so no admissible set contains \(c\).
Among subsets of \(\{a,b,d\}\), conflict-freeness excludes having both
\(a,b\). The set \(\{d\}\) cannot defend itself against \(c\),
whereas either \(a\) or \(b\) attacks \(c\) and also defends itself
against the other. Therefore the admissible sets are exactly
\[
 \varnothing,\quad\{a\},\quad\{b\},\quad\{a,d\},\quad\{b,d\}.
\]
Every argument has an attacker, so the empty set is fixed and grounded.
The characteristic values of \(\{a\}\) and \(\{a,d\}\) are both
\(\{a,d\}\): \(a\) defeats the attacker of \(a\) and that of
\(d\), while it cannot defend \(b\) or \(c\) against \(a\).
The symmetric calculation gives value \(\{b,d\}\) for the other two
nonempty admissible sets. Thus the complete sets are exactly the three
listed, and the two maximal sets are preferred. Each of those sets attacks
both outsiders, so both are stable. Their intersection is \(\{d\}\),
strictly larger than the empty grounded extension.
\end{proof}
\end{annexitem}
